\chapter*{Abstract}
\noindent\textbf{Title:} HDO VideoAPI: A Provider-Based Middleware for
Unified Access to Multiple Video Engines in Norwegian Emergency Medicine

\bigskip

\gls{hdo} operates video solutions that serve as decision support for \gls{amk}
and \gls{lvs} operators in Norway. The solutions are built on two video engines,
AntMedia Server and \gls{nhn} Video, which cover the same need but have different
\gls{api}s, authentication mechanisms, and domain models. Parallel integration
with both engines increases the complexity of operations and further development,
and raises the barrier to switching providers or adding new video engines.

This thesis investigates whether a \gls{provider}-based \gls{mellomvare} can give
HDO a single unified API for managing video rooms, while logic specific to each
video engine is kept separate from the rest of the system. The project was carried
out as a design and implementation project. The solution was developed in ASP.NET
Core with dedicated providers for AntMedia Server and NHN Video, \gls{token}-based
authentication, Docker support, health checks, structured logging, and Prometheus
metrics.

The evaluation consisted of requirements tracing, automated unit and integration
tests, end-to-end testing against HDO's test environments, and a final evaluation
from HDO. The architecture covers the core operations of creating, retrieving, and
deleting video rooms across both engines, and the provider layer isolates most of
the backend-specific logic. Full standardization is nevertheless difficult when the
video engines are built on different domain models. This is particularly true for
starting and stopping video streams, which do not carry the same meaning in
AntMedia and NHN Video.

A provider-based middleware reduces the coupling between HDO's client systems and
the underlying video engines, and makes it easier to support new video platforms in
the future. The result is a proof-of-concept demonstrating that the architecture is
feasible in practice. For the solution to be taken into production, it must be
further developed with stricter authorization, robust token validation, scalability
testing, and improved error handling against external video engines.